\section{Exhibit B-2: passivity toy (positivity ledger and zero confinement)}
\label{sec:exhibit-passivity}

A recurring \SBT motif is that the most rigidifying feasibility conditions are often \emph{positivity} conditions.
In physics, passivity/positivity constraints strongly restrict the analytic structure of response functions \cite{Willems1972II}.
In mathematics, many ``stability'' and ``no-gain'' requirements appear as nonnegativity of quadratic forms, positive-definite kernels, or positivity of coefficients in generating functions.
This exhibit provides a deliberately simple toy instantiation of that motif:
as a positivity constraint is tightened, zeros (defects) are increasingly confined to a symmetry locus.

\subsection{Self-duality and a symmetry locus}
We consider a family of polynomials
\[
Z_\lambda(z) \;=\; \sum_{k=0}^{n} a_k(\lambda)\,z^k
\]
constructed from a finite Ising-style partition sum (implemented by direct enumeration for $n=\PassivityN$ in the recorded run).
A key structural feature is a built-in self-duality symmetry:
we enforce palindromic coefficients by symmetrization,
\[
a_k(\lambda)\leftarrow \tfrac12\bigl(a_k(\lambda)+a_{n-k}(\lambda)\bigr),
\]
which makes $Z_\lambda$ invariant (up to the trivial $z^n$ factor) under inversion $z\leftrightarrow 1/z$.
Consequently, zeros occur in reciprocal pairs.
The fixed locus of this symmetry is the unit circle $|z|=1$, which we treat as the \emph{symmetry axis} (or symmetry manifold) for the toy.

In \SBT language, inversion symmetry is a \Ptwo\ constraint, while palindromization is a \Pone\ operator rewrite that makes the constraint manifest at the packaged level.

\subsection{Positivity as a feasibility ledger}
The toy feasibility constraint is a one-parameter ``positivity tightening'' of couplings.
We sample a set of edge couplings $J_e$ (both signs), and define a clipped family
\[
J_e(\lambda) \;=\; \max(J_e,0) + \lambda\,\min(J_e,0),
\qquad \lambda\in[0,1].
\]
Thus $\lambda=1$ is the original mixed-sign instance, while $\lambda=0$ is a purely nonnegative (ferromagnetic) instance.
This is a minimal model of a passivity/positivity ledger: tightening $\lambda\downarrow 0$ removes ``gain-like'' sign structure.

\subsection{Defects as zeros and a confinement metric}
We compute all complex zeros of $Z_\lambda$ and quantify confinement to the symmetry locus by the radial deviation
\[
d_i(\lambda) := \bigl||z_i(\lambda)|-1\bigr|,
\]
summarized by
\[
\mathrm{mean\_dev}(\lambda)=\frac1n\sum_i d_i(\lambda),
\qquad
\mathrm{max\_dev}(\lambda)=\max_i d_i(\lambda).
\]
These play the role of \Psix\ accounting diagnostics: they measure how strongly the defect set (zeros) respects the symmetry locus.

Table~\ref{tab:passivity} reports the confinement metrics across $\lambda$.
In the recorded run, tightening positivity produces a dramatic confinement effect:
\[
\mathrm{mean\_dev}(1) \approx \PassivityMeanDevLamOne,
\qquad
\mathrm{mean\_dev}(0) \approx \PassivityMeanDevLamZero,
\]
so the feasible ($\lambda\to 0$) regime concentrates zeros extremely close to $|z|=1$.

% Auto-generated by scripts/export_results_tex.py from notes/*_last_run.json
\begin{table}[H]
\centering
\caption{Passivity toy zero-confinement metrics}
\label{tab:passivity}
\begin{tabular}{l|l|l}
\hline
 $\lambda$ & mean\_dev & max\_dev \\
\hline
 \ensuremath{1} & \ensuremath{0.871694} & \ensuremath{7.10509} \\
\hline
 \ensuremath{0.7} & \ensuremath{0.580068} & \ensuremath{4.26786} \\
\hline
 \ensuremath{0.4} & \ensuremath{0.349525} & \ensuremath{2.26657} \\
\hline
 \ensuremath{0.2} & \ensuremath{0.200513} & \ensuremath{1.21606} \\
\hline
 \ensuremath{0.1} & \ensuremath{0.0946892} & \ensuremath{0.718256} \\
\hline
 0 & \ensuremath{3.0292\times 10^{-4}} & \ensuremath{9.0919\times 10^{-4}} \\
\hline
\end{tabular}
\end{table}


\begin{figure}[H]
\centering
\includegraphics[width=0.75\linewidth]{passivity_toy_last_run.png}%
\caption{Passivity toy: zeros of a self-dual polynomial family under positivity tightening. The symmetry locus is the unit circle $|z|=1$. As the clipping parameter $\lambda$ decreases toward the nonnegative (ferromagnetic) regime, the zeros concentrate near the symmetry locus; Table~\ref{tab:passivity} quantifies this via radial deviation metrics.}
\label{fig:passivity-toy}
\end{figure}

\subsection{Interpretation and scope}
This toy is not a theorem and does not claim to reproduce the full hypotheses of classical zero-confinement results.
Its role is methodological: it demonstrates a recognizable \SBT pattern in a controlled setting.
This motif has rigorous instances, most famously the Lee--Yang circle theorem for ferromagnetic Ising models \cite{LeeYang1952}; we do not invoke that theorem here.

\begin{itemize}
  \item Symmetry (\Ptwo) specifies a distinguished locus (here $|z|=1$).
  \item Positivity/passivity acts as a feasibility constraint measured by an audit ledger (\Psix).
  \item Under tighter feasibility, defects (zeros) become increasingly confined to the symmetry locus.
\end{itemize}

The purpose of including this exhibit in a mathematics instantiation paper is to motivate (without overclaiming) why problems like ``zero confinement to a self-dual manifold'' naturally invite passivity/positivity formulations.
Exhibit~B-1 on prime closures shows that naive staging and packaging can fail catastrophically in the critical strip; this toy suggests what a successful closure might require in addition: a feasibility ledger strong enough to prevent defect drift away from symmetry.
